\documentclass[12pt, a4paper]{article}
% Shared preamble — included by all Durin's Code LaTeX documents
\usepackage[a4paper, margin=2.5cm]{geometry}
\usepackage{amsmath, amssymb}
\usepackage{graphicx}
\usepackage{xcolor}
\usepackage{listings}
\usepackage{booktabs}
\usepackage{longtable}
\usepackage{array}
\usepackage{enumitem}
\usepackage{titlesec}
\usepackage{fancyhdr}
\usepackage{hyperref}
\usepackage{fontenc}
\usepackage{inputenc}
\usepackage{microtype}
\usepackage{caption}
\usepackage{subcaption}
\usepackage{tikz}
\usepackage{forest}
\usepackage{mdframed}
\usetikzlibrary{arrows.meta, positioning, shapes.geometric, fit, backgrounds}

% ── Colour palette ──────────────────────────────────────────────
\definecolor{durins-blue}{RGB}{31, 78, 121}
\definecolor{durins-gold}{RGB}{180, 130, 20}
\definecolor{durins-gray}{RGB}{90, 90, 90}
\definecolor{code-bg}{RGB}{248, 248, 248}
\definecolor{code-kw}{RGB}{0, 0, 180}
\definecolor{code-str}{RGB}{163, 21, 21}
\definecolor{code-cmt}{RGB}{0, 128, 0}
\definecolor{code-num}{RGB}{9, 134, 88}

% ── Listings style ──────────────────────────────────────────────
\lstdefinelanguage{DurinsCode}{
  morekeywords={room, action, item, npc, exit, if, else, print, remove,
                player, true, false, description, current_room},
  morestring=[b]",
  morecomment=[l]{//},
  sensitive=true
}
\lstdefinestyle{dc}{
  language=DurinsCode,
  backgroundcolor=\color{code-bg},
  basicstyle=\ttfamily\small,
  keywordstyle=\color{code-kw}\bfseries,
  stringstyle=\color{code-str},
  commentstyle=\color{code-cmt}\itshape,
  numberstyle=\tiny\color{durins-gray},
  numbers=left,
  numbersep=8pt,
  breaklines=true,
  frame=single,
  framerule=0.5pt,
  rulecolor=\color{durins-blue!40},
  xleftmargin=1.5em,
  aboveskip=0.8em,
  belowskip=0.8em,
  tabsize=4,
  showstringspaces=false,
  captionpos=b,
}
\lstdefinestyle{cpp}{
  language=C++,
  backgroundcolor=\color{code-bg},
  basicstyle=\ttfamily\footnotesize,
  keywordstyle=\color{code-kw}\bfseries,
  stringstyle=\color{code-str},
  commentstyle=\color{code-cmt}\itshape,
  numberstyle=\tiny\color{durins-gray},
  numbers=left, numbersep=8pt,
  breaklines=true,
  frame=single, framerule=0.5pt,
  rulecolor=\color{durins-blue!40},
  xleftmargin=1.5em,
  aboveskip=0.8em, belowskip=0.8em,
  tabsize=4, showstringspaces=false,
  captionpos=b,
}
\lstdefinestyle{bash}{
  language=bash,
  backgroundcolor=\color{code-bg},
  basicstyle=\ttfamily\small,
  keywordstyle=\color{code-kw},
  breaklines=true,
  frame=single, framerule=0.5pt,
  rulecolor=\color{durins-blue!40},
  xleftmargin=1.5em,
  numbers=none,
  aboveskip=0.8em, belowskip=0.8em,
  showstringspaces=false,
  captionpos=b,
}

% ── Section formatting ──────────────────────────────────────────
\titleformat{\section}
  {\color{durins-blue}\Large\bfseries}{{\thesection}}{1em}{}[\color{durins-blue}\titlerule]
\titleformat{\subsection}
  {\color{durins-blue!80}\large\bfseries}{{\thesubsection}}{1em}{}
\titleformat{\subsubsection}
  {\color{durins-gray}\normalsize\bfseries}{{\thesubsubsection}}{1em}{}

% ── Header / footer ─────────────────────────────────────────────
\pagestyle{fancy}
\fancyhf{}
\fancyhead[L]{\small\color{durins-gray}Durin's Code — CS4031 Compiler Construction}
\fancyhead[R]{\small\color{durins-gray}\leftmark}
\fancyfoot[C]{\small\color{durins-gray}\thepage}
\renewcommand{\headrulewidth}{0.4pt}
\renewcommand{\footrulewidth}{0.4pt}

% ── Hyperlink style ─────────────────────────────────────────────
\hypersetup{
  colorlinks=true,
  linkcolor=durins-blue,
  urlcolor=durins-blue,
  citecolor=durins-gold,
  pdftitle={Durin's Code — CS4031},
  pdfauthor={Huzaifa Abdul Rehman, Abdul Moiz Hussain, Muhammad Abdullah Khan},
}

% ── Handy boxes ────────────────────────────────────────────────
\newmdenv[
  linecolor=durins-blue!60,
  backgroundcolor=durins-blue!5,
  roundcorner=4pt,
  leftmargin=0, rightmargin=0,
  innerleftmargin=10pt, innerrightmargin=10pt,
  innertopmargin=8pt, innerbottommargin=8pt,
]{infobox}


\begin{document}

% ════════════════════════════════════════════════════════════════
%  TITLE PAGE
% ════════════════════════════════════════════════════════════════
\begin{titlepage}
  \centering
  \vspace*{2cm}
  {\color{durins-blue}\rule{\linewidth}{2pt}}\\[1em]
  {\Huge\bfseries\color{durins-blue} Durin's Code}\\[0.5em]
  {\LARGE\color{durins-gold} Language Reference Manual}\\[1em]
  {\color{durins-blue}\rule{\linewidth}{2pt}}\\[2em]

  {\large CS4031 --- Compiler Construction \\ Spring 2026}\\[3em]

  \begin{tabular}{ll}
    \toprule
    \textbf{Student Name} & \textbf{Student ID} \\
    \midrule
    Huzaifa Abdul Rehman  & 23K-0782 \\
    Abdul Moiz Hussain    & 23K-0553 \\
    Muhammad Abdullah Khan & 23K-0607 \\
    \bottomrule
  \end{tabular}

  \vfill
  {\small\color{durins-gray} Document version 1.0 \quad|\quad \today}
\end{titlepage}

\tableofcontents
\newpage

% ════════════════════════════════════════════════════════════════
\section{Introduction}
% ════════════════════════════════════════════════════════════════

Durin's Code (file extension \texttt{.dc}) is a domain-specific language designed for
authoring interactive text-adventure games in the spirit of the 8-bit era.
A \texttt{.dc} program declares a \emph{world} --- a directed graph of rooms populated with
items and NPCs --- and a set of \emph{actions}, which are named event handlers triggered by
player input at runtime.

The compiler performs full static validation of the world graph, generates optimised
Three-Address Code (TAC) as an intermediate representation, serialises the result to
JSON bytecode, and executes it inside a terminal-based Virtual Machine.

\begin{infobox}
\textbf{Design philosophy.}
Durin's Code removes the need for imperative low-level coding when building text
adventures. Every game element (rooms, items, NPCs, transitions) is expressed
declaratively. All cross-references are validated at compile time, preventing
runtime crashes from missing rooms or invalid items.
\end{infobox}

% ════════════════════════════════════════════════════════════════
\section{Lexical Structure}
% ════════════════════════════════════════════════════════════════

\subsection{Source Encoding}
Source files are plain text encoded in UTF-8.
Both Unix (\texttt{\textbackslash n}) and Windows (\texttt{\textbackslash r\textbackslash n})
line endings are accepted; the lexer normalises them to a single newline internally.

\subsection{Comments}
Only single-line comments are supported. A comment begins with \texttt{//} and
extends to the end of the current line.

\begin{lstlisting}[style=dc, caption={Comment syntax}]
// This is a comment. It is ignored by the compiler.
room "start" {   // inline comments are also allowed
    description "A quiet room."
}
\end{lstlisting}

\subsection{Identifiers}

\[
  \textit{identifier} \;\;::=\;\; [\_\,\text{a-zA-Z}]\;[\_\,\text{a-zA-Z0-9}]^*
\]

Identifiers are case-sensitive. A leading underscore is permitted. Examples:
\texttt{the\_ring}, \texttt{\_private}, \texttt{room1}, \texttt{hasItem}.

\subsection{Keywords}

The following identifiers are reserved and may not be used as user-defined names:

\begin{center}
\begin{tabular}{llll}
\toprule
\texttt{room} & \texttt{action} & \texttt{item} & \texttt{npc} \\
\texttt{exit} & \texttt{if} & \texttt{else} & \texttt{print} \\
\texttt{remove} & \texttt{player} & \texttt{true} & \texttt{false} \\
\texttt{description} & \texttt{current\_room} & & \\
\bottomrule
\end{tabular}
\end{center}

\subsection{Literals}

\paragraph{String literals.}
A string is enclosed in double-quotes on a single line.
\[
  \textit{STRING} \;\;::=\;\; \texttt{"} \; [^\texttt{"\textbackslash n}]^* \; \texttt{"}
\]

\paragraph{Integer literals.}
A sequence of decimal digits. Floating-point numbers are \textbf{not} supported and
produce a compile-time error.
\[
  \textit{NUMBER} \;\;::=\;\; [0\text{-}9]^+
\]

\paragraph{Boolean literals.} \texttt{true} and \texttt{false}.

\subsection{Operators and Punctuation}

\begin{center}
\begin{tabular}{cll}
\toprule
\textbf{Symbol} & \textbf{Name} & \textbf{Usage} \\
\midrule
\texttt{\{} \texttt{\}} & Braces & Block delimiters \\
\texttt{(} \texttt{)} & Parentheses & Function-call arguments \\
\texttt{,} & Comma & Property separator \\
\texttt{:} & Colon & Property value separator \\
\texttt{;} & Semicolon & Optional statement terminator \\
\texttt{.} & Dot & Attribute access (\texttt{player.inventory}) \\
\texttt{=} & Assign & Simple assignment \\
\texttt{==} & Equal & Equality test \\
\texttt{+=} & Plus-assign & Add to attribute / add item to inventory \\
\texttt{-=} & Minus-assign & Subtract from attribute / remove from inventory \\
\texttt{+} \texttt{-} & Arithmetic & Integer arithmetic \\
\texttt{>} \texttt{>=} \texttt{<} \texttt{<=} & Relational & Numeric comparisons \\
\texttt{\&\&} & Logical AND & Both conditions must hold \\
\texttt{||} & Logical OR & Either condition must hold \\
\bottomrule
\end{tabular}
\end{center}

\subsection{Context-Sensitive Dot Tokens}

After a \texttt{.} the lexer inspects the immediately preceding token to disambiguate
attribute access:

\begin{itemize}
  \item \texttt{player.}\textit{attr} $\Rightarrow$ next identifier emitted as \texttt{TOKEN\_PLAYER\_ATTR}
  \item \textit{room\_id}\texttt{.}\textit{attr} $\Rightarrow$ next identifier emitted as \texttt{TOKEN\_ROOM\_ATTR}
\end{itemize}

This allows the parser to differentiate \texttt{player.has\_item(\ldots)} from
\texttt{room.description} without extra lookahead.

% ════════════════════════════════════════════════════════════════
\section{Grammar (EBNF)}
% ════════════════════════════════════════════════════════════════

The complete context-free grammar of Durin's Code in Extended Backus-Naur Form:

\begin{lstlisting}[language={}, basicstyle=\ttfamily\small, frame=single,
    backgroundcolor=\color{code-bg}, caption={Complete EBNF Grammar}]
program         ::=  declaration* EOF

declaration     ::=  room_decl | action_decl

(* World *)
room_decl       ::=  "room" STRING "{" room_body* "}"
room_body       ::=  description_stmt
                   | item_decl
                   | npc_decl
                   | exit_decl

description_stmt ::= "description" STRING

item_decl       ::=  "item" IDENTIFIER "{" prop_list "}"
npc_decl        ::=  "npc"  IDENTIFIER "{" prop_list "}"
prop_list       ::=  prop ("," prop)*
prop            ::=  IDENTIFIER ":" (NUMBER | STRING | BOOLEAN)

exit_decl       ::=  "exit" IDENTIFIER STRING
                (*          ^direction  ^target_room_id *)

(* Actions *)
action_decl     ::=  "action" STRING "{" statement* "}"

statement       ::=  print_stmt
                   | remove_stmt
                   | assign_stmt
                   | if_stmt

print_stmt      ::=  "print" STRING
remove_stmt     ::=  "remove" IDENTIFIER
assign_stmt     ::=  player_attr_ref ("=" | "+=" | "-=")
                     (IDENTIFIER | NUMBER | STRING | BOOLEAN)
player_attr_ref ::=  "player" "." IDENTIFIER

if_stmt         ::=  "if" condition "{" statement* "}"
                     ( "else" "{" statement* "}" )?

(* Conditions *)
condition       ::=  simple_cond ( ("&&" | "||") simple_cond )*
simple_cond     ::=  room_check | has_item_check | comparison

room_check      ::=  "current_room" "==" STRING
has_item_check  ::=  "player" "." "has_item" "(" IDENTIFIER ")"
comparison      ::=  player_attr_ref
                     ("==" | ">" | ">=" | "<" | "<=")
                     (NUMBER | STRING | BOOLEAN)

(* Terminals *)
STRING          ::=  '"' [^"\n]* '"'
NUMBER          ::=  [0-9]+
BOOLEAN         ::=  "true" | "false"
IDENTIFIER      ::=  [_a-zA-Z] [_a-zA-Z0-9]*
\end{lstlisting}

% ════════════════════════════════════════════════════════════════
\section{Declarations}
% ════════════════════════════════════════════════════════════════

\subsection{Room Declaration}

A room declaration defines one node in the world graph.

\begin{lstlisting}[style=dc, caption={Room declaration syntax}]
room "bag_end" {
    description "The cozy hole of a Hobbit. A golden ring glints on the table."
    item the_ring { power: 100, type: "artifact" }
    npc gandalf   { health: 500, hostile: false }
    exit east  "buckland"
    exit north "shire"
}
\end{lstlisting}

\begin{itemize}
  \item \textbf{\texttt{description}} — Sets the text printed by the \texttt{look} command.
        Exactly one description per room.
  \item \textbf{\texttt{item}} — Places a named item in this room with a property map.
        The item is globally registered and may be referenced in any action.
  \item \textbf{\texttt{npc}} — Places a named NPC in the room with a property map.
  \item \textbf{\texttt{exit}} — Declares a directed edge; direction is a free identifier
        (\texttt{north}, \texttt{east}, or any word).
\end{itemize}

\paragraph{Semantic constraints.}
\begin{enumerate}
  \item Room identifiers (the string name) must be unique across the program.
  \item Exit target strings must reference a declared room.
  \item No two exits in the same room may share the same direction identifier.
\end{enumerate}

\subsection{Action Declaration}

An action is a named event handler triggered by matching player input.

\begin{lstlisting}[style=dc, caption={Action declaration syntax}]
action "take ring" {
    if current_room == "bag_end" {
        player.inventory += the_ring
        print "The Precious is yours."
    }
}
\end{lstlisting}

Action names are strings and may contain spaces. At runtime the player types the action
name verbatim (case-sensitive).

% ════════════════════════════════════════════════════════════════
\section{Statements}
% ════════════════════════════════════════════════════════════════

\subsection{\texttt{print}}
Writes a string literal to standard output.
\begin{lstlisting}[style=dc, numbers=none]
print "You cannot do that here."
\end{lstlisting}

\subsection{\texttt{remove}}
Removes a named item from the game world and from the player's inventory if present.
\begin{lstlisting}[style=dc, numbers=none]
remove the_ring
\end{lstlisting}
\textbf{Semantic constraint:} the identifier must name a declared item.

\subsection{Assignment}
Sets or modifies a player attribute.

\begin{lstlisting}[style=dc, numbers=none, caption={Assignment variants}]
player.win       = true          // boolean assignment
player.health    = 100           // integer assignment
player.name      = "Bilbo"       // string assignment
player.inventory += the_ring     // add item to inventory
player.health    -= 10           // subtract from integer attribute
\end{lstlisting}

\subsection{\texttt{if} / \texttt{else}}
Conditional execution. The \texttt{else} branch is optional.

\begin{lstlisting}[style=dc, caption={Conditional statement}]
if current_room == "mount_doom" && player.has_item(the_ring) {
    remove the_ring
    player.win = true
    print "Middle-earth is saved!"
} else {
    print "You cannot do that here."
}
\end{lstlisting}

% ════════════════════════════════════════════════════════════════
\section{Conditions}
% ════════════════════════════════════════════════════════════════

\begin{center}
\begin{tabular}{p{5.5cm} p{7cm}}
\toprule
\textbf{Form} & \textbf{Meaning} \\
\midrule
\texttt{current\_room == "id"} & Player is in room \texttt{"id"} \\
\texttt{player.has\_item(name)} & Player's inventory contains \texttt{name} \\
\texttt{player.attr == value} & Attribute equality \\
\texttt{player.attr > value} & Greater-than comparison \\
\texttt{player.attr >= value} & Greater-or-equal \\
\texttt{player.attr < value} & Less-than comparison \\
\texttt{player.attr <= value} & Less-or-equal \\
\texttt{cond1 \&\& cond2} & Logical AND \\
\texttt{cond1 || cond2} & Logical OR \\
\bottomrule
\end{tabular}
\end{center}

% ════════════════════════════════════════════════════════════════
\section{Type System}
% ════════════════════════════════════════════════════════════════

Property values in \texttt{item} and \texttt{npc} blocks must be one of three scalar types:

\begin{center}
\begin{tabular}{lll}
\toprule
\textbf{Type} & \textbf{Example} & \textbf{Storage} \\
\midrule
Integer & \texttt{100} & 32-bit signed \\
String & \texttt{"artifact"} & \texttt{std::string} \\
Boolean & \texttt{true} / \texttt{false} & \texttt{bool} \\
\bottomrule
\end{tabular}
\end{center}

Player attributes are dynamically typed: the first assignment determines the type.
Floating-point numbers are rejected at lex time.

% ════════════════════════════════════════════════════════════════
\section{Scope Rules}
% ════════════════════════════════════════════════════════════════

\begin{itemize}
  \item All room, item, NPC, and action names share a single flat global namespace.
  \item Duplicate names within the same category are a semantic error.
  \item Items declared inside any room are globally visible in all actions.
  \item \texttt{player} is a built-in identifier; \texttt{player.X} accesses runtime state.
  \item There is no block-level scoping — no local variables exist in actions.
\end{itemize}

% ════════════════════════════════════════════════════════════════
\section{Error Catalogue}
% ════════════════════════════════════════════════════════════════

\subsection{Lexer Errors}

\begin{center}
\begin{tabular}{p{5cm} p{8cm}}
\toprule
\textbf{Message} & \textbf{Cause} \\
\midrule
\texttt{Unexpected character} & A character outside the language alphabet (e.g.\ \texttt{@}) \\
\texttt{Unterminated string.} & A \texttt{"} opened but not closed before end-of-line \\
\texttt{Floating-point numbers are not supported.} & A digit sequence containing \texttt{.} \\
\texttt{Expected '\&\&'} & Single \texttt{\&} not followed by \texttt{\&} \\
\texttt{Expected '||'} & Single \texttt{|} not followed by \texttt{|} \\
\bottomrule
\end{tabular}
\end{center}

\subsection{Semantic Errors}

\begin{center}
\begin{tabular}{p{5.5cm} p{7.5cm}}
\toprule
\textbf{Error} & \textbf{Cause} \\
\midrule
Duplicate room & Two rooms share the same name \\
Duplicate item & Two items share the same name \\
Duplicate action & Two actions share the same name \\
Exit target not declared & An exit points to a non-existent room \\
Duplicate exit direction & A direction is used twice in the same room \\
\texttt{remove}: unknown item & Removal of an undeclared item \\
\texttt{has\_item}: unknown item & Inventory check against an undeclared item \\
\texttt{current\_room}: unknown room & Room comparison against an undeclared room \\
\bottomrule
\end{tabular}
\end{center}

% ════════════════════════════════════════════════════════════════
\section{Complete Example Program}
% ════════════════════════════════════════════════════════════════

\begin{lstlisting}[style=dc, caption={Middle-Earth Adventure --- full demo program}]
// Middle-earth adventure

room "bag_end" {
    description "The cozy hole of a Hobbit. A golden ring glints on the table."
    item the_ring { power: 100, type: "artifact" }
    exit east "buckland"
}

room "buckland" {
    description "A quiet riverside village. The old ferry creaks at the dock."
    exit west "bag_end"
    exit east "mount_doom"
}

room "mount_doom" {
    description "The air is thick with ash. The fiery Cracks of Doom loom ahead."
    npc sauron { health: 999, hostile: true }
    exit west "buckland"
}

action "take ring" {
    if current_room == "bag_end" {
        player.inventory += the_ring
        print "The Precious is yours."
    } else {
        print "There is nothing of value here."
    }
}

action "destroy ring" {
    if current_room == "mount_doom" && player.has_item(the_ring) {
        remove the_ring
        print "The ring is consumed by fire! Middle-earth is saved."
        player.win = true
    } else {
        print "You cannot do that here."
    }
}
\end{lstlisting}

\end{document}
